\chapter{リアクションホイールと AMBAC の統一的理解（Sim03）}\label{ch:rw}

\section{同じ式、違う関節}

導出は付録~\ref{app:rw}（ホイールは曲率ゼロの関節）と付録~\ref{app:saturation}（角運動量の容量）を参照。

\begin{figure}[htbp]\centering
\resizebox{0.8\linewidth}{!}{% リアクションホイールと手足
\begin{tikzpicture}[>=Stealth,font=\small]
\begin{scope}
  \draw[thick,fill=gray!15] (-0.5,-1) rectangle (0.5,1);
  \draw[thick,fill=orange!30] (0,0) circle (0.35);
  \draw[->,thick] (0.5,0) arc (0:300:0.5);
  \fill (0,0) circle (1.5pt);
  \node[align=center,font=\footnotesize,below] at (0,-1.15) {リアクションホイール\\重心が回転軸上：形が変わらない\\$\Gamma=$ 一定、曲率 $F=0$\\何回転でも回せる};
\end{scope}
\begin{scope}[shift={(5,0)}]
  \draw[thick,fill=gray!15] (-0.5,-1) rectangle (0.5,1);
  \fill (0,0.6) circle (2pt);
  \draw[line width=3pt,blue!60] (0,0.6)--++(20:1.3);
  \draw[line width=3pt,blue!25] (0,0.6)--++(80:1.1);
  \draw[->,thick] (0,0.6)++(20:0.9) arc (20:80:0.9);
  \node[align=center,font=\footnotesize,below] at (0,-1.15) {手足（AMBAC）\\重心が動く：形が変わる\\$\Gamma(q)$ が形に依存、$F\neq0$\\可動域がある};
\end{scope}
\node[font=\footnotesize] at (2.5,0.4) {$\dot\theta=-\Gamma\,\dot q$};
\end{tikzpicture}
}
\caption{リアクションホイールと手足（AMBAC）。どちらも $\dot\theta=-\Gamma\dot q$ で書けるが、関節の性質が異なる。}\label{fig:tikz-rw}
\end{figure}


人工衛星のリアクションホイール（RW）と AMBAC は、どちらも \textbf{運動量交換} による姿勢制御で、同じ式で書ける。

\[
\dot\theta = \frac{L - B(q)\cdot\dot q}{A(q)} = \frac{L}{A(q)} - \Gamma(q)\cdot\dot q
\]

\begin{longtable}[]{@{}
  >{\raggedright\arraybackslash}p{(\columnwidth - 4\tabcolsep) * \real{0.3333}}
  >{\raggedright\arraybackslash}p{(\columnwidth - 4\tabcolsep) * \real{0.3333}}
  >{\raggedright\arraybackslash}p{(\columnwidth - 4\tabcolsep) * \real{0.3333}}@{}}
\toprule\noalign{}
\begin{minipage}[b]{\linewidth}\raggedright
\end{minipage} & \begin{minipage}[b]{\linewidth}\raggedright
リアクションホイール
\end{minipage} & \begin{minipage}[b]{\linewidth}\raggedright
AMBAC（四肢）
\end{minipage} \\
\midrule\noalign{}
\endhead
\bottomrule\noalign{}
\endlastfoot
関節 & 連続回転できる & 可動域がある \\
回したときの質量配置 & 変わらない（重心が回転軸上） & 変わる \\
接続 Γ & \textbf{一定}（I\_w / A） & \textbf{形状 q に依存} \\
曲率 F & 0 & ≠ 0 \\
任意角の回転 & 回し続ければ可能（ホロノミック） & ループ運動で可能（非ホロノミック・幾何学的位相） \\
たとえ & 低ギア・無限ストローク & 高ギア・短ストローク \\
\end{longtable}

\textbf{RW は「AMBAC の特殊な場合」}、すなわち曲率ゼロ・可動域無限の関節とみなせる。 逆に AMBAC は「可動域を補うために、形状空間のループで回転を稼ぐ RW」と言える。

\section{Part A: 接続とループ最適化}

\FIG{sim03_connection.png}

関節可動域は腕 −150°〜150°、脚 −120°〜30°。

\begin{longtable}[]{@{}
  >{\raggedright\arraybackslash}p{(\columnwidth - 2\tabcolsep) * \real{0.5000}}
  >{\raggedright\arraybackslash}p{(\columnwidth - 2\tabcolsep) * \real{0.5000}}@{}}
\toprule\noalign{}
\begin{minipage}[b]{\linewidth}\raggedright
項目
\end{minipage} & \begin{minipage}[b]{\linewidth}\raggedright
値
\end{minipage} \\
\midrule\noalign{}
\endhead
\bottomrule\noalign{}
\endlastfoot
RW の Γ\_w & 0.0014（ホイール 1 回転で胴体 0.51°） \\
脚の Γ\textsubscript{2} & 約 0.33（脚を 1° 振ると胴体 0.33°） \\
Sim02 の長方形ループ & +2.98°/周 \\
\textbf{最適円ループ}（中心 腕 −20°, 脚 −40°、半径 70°） & \textbf{+6.70°/周} \\
可動域内の理論上界（同符号曲率の全領域を囲む） & +9.57°/周 \\
\end{longtable}

\begin{itemize}
\tightlist
\item
  四肢の Γ は RW の 200 倍以上ある。四肢は「1 回の動きで大きく向きを変える」装置。
\item
  ループ最適化で 1 周あたりの回転は 2.25 倍になった。理論上界の 70\% に達している。 曲率マップの破線（F=0）が「囲むべき領域」の境界になっている。
\item
  上界は関節可動域で決まる。\textbf{AMBAC 性能は可動域設計の問題}でもある。
\end{itemize}

\section{Part B: 90° の姿勢変更（静止→静止）}

\FIG{sim03_rest2rest.png}

\begin{longtable}[]{@{}
  >{\raggedright\arraybackslash}p{(\columnwidth - 4\tabcolsep) * \real{0.3333}}
  >{\raggedright\arraybackslash}p{(\columnwidth - 4\tabcolsep) * \real{0.3333}}
  >{\raggedright\arraybackslash}p{(\columnwidth - 4\tabcolsep) * \real{0.3333}}@{}}
\toprule\noalign{}
\begin{minipage}[b]{\linewidth}\raggedright
\end{minipage} & \begin{minipage}[b]{\linewidth}\raggedright
RW
\end{minipage} & \begin{minipage}[b]{\linewidth}\raggedright
AMBAC
\end{minipage} \\
\midrule\noalign{}
\endhead
\bottomrule\noalign{}
\endlastfoot
所要時間 & 10 s（設定値） & 26 s（13 周 × 0.5 Hz） \\
実行量 & ホイール 178 回転、最大 2003 rpm & 四肢ループ 13 周 \\
最終姿勢 & 90.00° & 87.16°（整数周回の端数） \\
途中の姿勢 & 単調・滑らか & \textbf{最大 113°（周回ごとに ±20〜25° 揺れる）} \\
\end{longtable}

\textbf{考察} - AMBAC では、正味の回転の上に、形状変化による大きな \textbf{可逆的な揺れ} が重なる。 観測・照準・作業を伴う姿勢制御では、この揺れは実用上の大問題になる。 対策としては、ループ形状を「揺れが小さく曲率が大きい」ものに最適化すること、 他の肢で揺れを打ち消すこと（冗長肢の利用）が考えられる。 - 端数（87° vs 90°）は、最後の部分ループや単関節の可逆運動で微調整できる （可逆 AMBAC＝微調整、非ホロノミック AMBAC＝大角度、という役割分担）。

\section{Part C: 外乱トルク下の姿勢保持と「飽和」}

\FIG{sim03_disturbance.png}

一定の外乱トルク τ\_d = 0.01 N·m を胴体に加える。全角運動量は L(t) = τ\_d t で増え続け、 \textbf{内部運動では L を消せない}。姿勢を保つには、増えた L をどこかに「溜める」しかない。

\begin{itemize}
\tightlist
\item
  \textbf{RW}: ホイールの回転に溜める。ホイール回転数が時間に比例して増え、ω\_max に達すると飽和する。
\item
  \textbf{AMBAC}: \textbf{四肢の周期運動に溜める}。ループ周回周波数 f が時間に比例して増え、f\_max に達すると飽和する。
\end{itemize}

AMBAC の平均化モデル（1 周平均で dθ̄/dt ≈ f·Δθ\_c + L·$\langle$1/A$\rangle$）から、角運動量容量は次のとおり。

\[
h_\mathrm{RW} = I_w\,\omega_{\max}, \qquad
h_\mathrm{AMBAC} = \frac{f_{\max}\,|\Delta\theta_c|}{\langle 1/A\rangle}
\]

\begin{longtable}[]{@{}
  >{\raggedright\arraybackslash}p{(\columnwidth - 4\tabcolsep) * \real{0.3333}}
  >{\raggedright\arraybackslash}p{(\columnwidth - 4\tabcolsep) * \real{0.3333}}
  >{\raggedright\arraybackslash}p{(\columnwidth - 4\tabcolsep) * \real{0.3333}}@{}}
\toprule\noalign{}
\begin{minipage}[b]{\linewidth}\raggedright
\end{minipage} & \begin{minipage}[b]{\linewidth}\raggedright
角運動量容量
\end{minipage} & \begin{minipage}[b]{\linewidth}\raggedright
0.01 N·m での保持限界
\end{minipage} \\
\midrule\noalign{}
\endhead
\bottomrule\noalign{}
\endlastfoot
RW（2 kg、5700 rpm） & 8.64 kg·m²/s & 864 s \\
AMBAC（腕 4 kg + 脚 12 kg、1 Hz） & 1.18 kg·m²/s & 118 s \\
\end{longtable}

保持中の AMBAC は、周回平均姿勢 \textbar θ̄\textbar{} ≤ 0.30° に保てたが、瞬時姿勢は ±23° 揺れる。

\textbf{考察} - 「AMBAC で回転を止められない」ことと「RW の運動量飽和」は \textbf{同一の物理現象} である。 RW が角運動量をホイールの回転に溜めるのと同じく、AMBAC は四肢の周期運動に溜める （連続的なループ運動は、実効的に「回り続けるホイール」になっている）。 - この人間大モデルでは、2 kg の RW の方が 16 kg の四肢より 7 倍以上の容量を持つ。 AMBAC の強みは容量ではなく、\textbf{推進剤も専用機器も使わず、既存の四肢で姿勢を変えられる}ことにある。 - どちらも飽和したら、外部トルク（スラスタ、磁気トルカ、接触）による \textbf{アンローディング} が必要になる。 劇中で AMBAC とスラスタが併用されていることは、物理的に自然な設計と言える。 - 周回内の揺れの補償には、ループ形状から計算した周期成分 θ\_osc(φ) を使った。 \textbf{幾何構造（接続・曲率）を知っていると、制御設計が格段に簡単になる}。 これはフィジカル AI（学習で方策を得る場合）が「何を内部表現として獲得すべきか」という問いに直結する。

\section{次への論点}

\begin{enumerate}
\def\labelenumi{\arabic{enumi}.}
\tightlist
\item
  \textbf{トルクとエネルギー（Sim04）}: f\_max は本来アクチュエータのトルク・出力で決まる。 順動力学で、関節トルク・消費エネルギー・f\_max を見積もる。
\item
  \textbf{スケーリング（18 m 級）}: Δθ\_c はスケール不変だが、A ∝ M ℓ² であり、f\_max もスケールに依存する。 巨大ロボで h\_AMBAC がどう伸びるかは、アクチュエータのスケーリング則次第（Sim04 で検証予定）。
\item
  \textbf{揺れの小さいループ}: 「正味回転は大きく、周回内の揺れは小さい」ループ形状の多目的最適化。
\item
  \textbf{フィジカル AI}: RW＋AMBAC＋スラスタの統合制御（運動量の配分）を、最適制御と学習でそれぞれ解いて比較する。
\end{enumerate}
